The evaluation combines archived checkpoints, known-signal calibration and independent ridge models, followed by construct controls and repeated-map sensitivity. Fixed-checkpoint evaluation reuses 45 aligned states from five Time-IMM datasets, three mechanisms and three training seeds. Eight conditions yield 360 evaluations: aligned input, the original restricted shuffle, four nominal derangement doses (25/50/75/100\%), zero content and a training-derived centroid. Numerical inputs, targets and event metadata remain fixed. Effect intervals use 2,000 paired bootstrap draws with entity resampling and within-entity circular window blocks of lengths 2, 4 and 8. Seeds are paired and averaged within a unit; their variation is reported separately. These exploratory intervals condition on the selected checkpoints, retained entities and realized maps. Archived validation losses against persistence and condition-retrained time-only models appear in Table~\ref{tab:historical_validation}. Formal capability screening remains \emph{Not assessed}; initialized selections are identified separately.

Known-signal calibration tests diagnostic sensitivity under specified truth. Eight EPA-Air event-time templates supply irregular numerical and text schedules. Numerical values, direction-bearing prose, targets and arrival delays are generated independently for each synthetic entity. Each of 20 independently seeded worlds contains 128 training, 32 validation and 128 test entities assigned across the fixed schedules. Four signal strengths $\gamma\in\{0,0.25,0.5,1\}$ share generated inputs within each world, permitting paired dose comparisons at fixed event geometry. Numerical targets are synthetic rather than observed environmental outcomes.

At a query origin, eligible directions are exponentially pooled with a one-day recency scale. The target is $0.8$ times the latest numerical value plus $0.2$ times its history mean, plus $\gamma$ times the variance-normalized eligible direction signal and independent Gaussian noise with standard deviation 0.35. Numerical histories follow an AR(1) process with coefficient 0.985. Independent simulated arrival delay has mean 0.15 days; queries cannot access an unarrived or future event. This generator makes text recoverability explicit while keeping random direction independent of timing and occurrence.

Oracle direction vectors and frozen MiniLM text representations are evaluated separately. Four paired increase/decrease phrasings are used for training and four different pairs for testing. PCA uses training phrases only, and standardized linear predictors are fitted by SGD with L2 regularization $\alpha=0.003$, at most 300 iterations and three fixed optimizer seeds. A numerical/timing comparator, an aligned predictor and a retrained shuffled predictor produce 1,200 fits across the full matrix. These fitting settings are fixed; validation data select neither hyperparameters nor checkpoints in this batch. The primary within-world test is a one-sided paired entity-level $t$ test at 0.05 after averaging optimizer-seed losses within each entity. Across-world detection rates use Wilson intervals. A frozen validation screen for incremental content benefit is applied separately, and both screened and unscreened detection are reported.

\paragraph{Validation screen for incremental content benefit.}
For each generated world, let $\ell^{\rm NT}_{V,e}$ and $\ell^A_{V,e}$ be the numerical/timing and aligned mean validation losses of entity $e$, with windows averaged within entities and optimizer seeds averaged before comparison where applicable. With $n_V=32$, the implemented rule is
\begin{equation}
 \begin{aligned}
 d_e&=\ell^{\rm NT}_{V,e}-\ell^A_{V,e},\qquad
 b_V=\bar d-t_{0.975,31}\frac{s_d}{\sqrt{32}},\\
 r_V&=\frac{\bar\ell^{\rm NT}_V-\bar\ell^A_V}{\bar\ell^{\rm NT}_V}.
 \end{aligned}
 \label{eq:validation_content_screen}
\end{equation}
where $s_d$ is the sample standard deviation of the paired entity differences. The screen requires $b_V>0$ and $r_V>0.01$. It concerns incremental content benefit over a comparator retaining numerical history and event timing. Overall predictive capability is a separate comparison with a task-appropriate forecasting baseline, here persistence for the real-data models. The 1\% margin is a batch-specific screening choice. Failure of either comparison leaves a fixed-model perturbation effect defined, but limits interpreting a small effect as evidence that a capable model ignored available content.

Independent residual-ridge models assess representation and dataset sensitivity, retaining all ten local FNSPID entities and eight EPA-Air entities. Each summary uses the latest four historical notes, with counts and recency retained separately. We compare cached GPT-2 with pinned MiniLM and FinBERT representations \cite{reimers2019sentencebert,minilmmodel,finbertmodel}. Modern encoders use attention-mask token mean pooling, L2 normalization, a 256-token limit and 32-dimensional training-only PCA. FinBERT's classifier head is excluded; its pooled hidden states are not sentence-trained embeddings. Penalties $1,10,100,1000$ are selected on pooled observed-scalar validation MSE. Condition-retrained models have their own validation selection; fixed perturbations retain the aligned checkpoint. Persistence, numerical-only ridge and numerical/timing ridge provide distinct comparator evidence. Their saved validation means address overall prediction and incremental content benefit; no paired validation screen was applied (Table~\ref{tab:ridge_capability_details}).

The same independent modeling route is applied to all nine Fidel-TS plants using numerical history and published weather embeddings. It uses 360 history hours to forecast the next 24 hours at daily origins, the original year splits, and reports whose recorded timestamp is no later than the numerical cutoff. This is a timestamp-restricted published-embedding experiment, not a reproduction of TGTSF or proof of report issuance time. Report validity timestamps are the available metadata; actual publication and arrival timestamps are absent. Shared weather context also precludes treating all plant windows as independent replications.

New configurations and script identities were recorded before their respective scoring batches. The local records are execution freezes, not externally registered studies. The Time-IMM and Fidel-TS historical test sets had already been inspected, whereas the controlled calibration worlds are newly generated. The distinction is retained in the result records.

\subsection{Construct-specific validation and repeated intervention maps}

Construct validation reuses the eight EPA-Air time templates with independent numerical innovations, per-note direction/source signs, sentence assignments, arrival delays and target noise. Twenty worlds each contain 128 training, 32 validation and 128 test entities, with nonoverlapping seven-day histories. Targets add one of three terms to the numerical baseline: variance-normalized recency-pooled direction, its independently generated source counterpart, or training-standardized log available-note count. Noise standard deviation is 0.35. Sixteen directional/source sentences train the representation pipeline, sixteen different directional phrasings provide validation/test inputs and sixteen further phrasings provide paraphrases. The source vocabulary deliberately stays fixed across splits.

Frozen MiniLM pooling, training-only PCA (at most fifteen components) and scaling precede ridge selection from $0.0001,0.01,1,100$ on mean validation entity loss. No latent direction/source readout enters the sentence predictor. The predictor and numerical/timing comparator each undergo validation selection, yielding 480 candidate fits and 120 selected states. The oracle readout directly evaluates the known generator signal without fitting.

The selected models are then screened with Eq.~\ref{eq:validation_content_screen} on the same 32 validation entities used to choose both ridge penalties. Consequently, $b_V$ is a \emph{post-selection nominal lower bound}. The screen is descriptive: its usual two-sided 95\% calculation does not provide an unadjusted 95\% coverage guarantee after this selection. The local freeze fixes the rule but does not remove selection effects. All independent test entities and all predefined control outcomes remain in the reported results, regardless of screening status.

Each aligned predictor evaluates whole permutation, direction flip, source flip, meaning-preserving paraphrase and zero at fixed event slots. In both generated batches, whole permutation is a uniform nonzero cyclic rotation within the current available history, retaining empty/singleton cases. Targeted edits preserve the alternative generated cue. Tests pair mean losses over 128 independently generated test entities within a world; effect intervals use twenty independent worlds, conditional on the same empirical clocks. All predefined cells are reported, with descriptive rejection rates and no universal multiplicity or error-control claim.

The map-sensitivity batch retains the independent aligned ridge states on EPA-Air, FNSPID and Fidel-TS and draws twenty additional maps with seeds 120001--120020. Each entity/split receives an independent uniform nonzero circular shift of ordered content; the same shifts are shared across representations. Fidel shifts the common report sequence within split and shares its map across all nine plants. These draws supplement the original map, without refitting or validation reselection. They report empirical 2.5th--97.5th percentiles and sign counts over maps, separately from population or training uncertainty. Donor timestamps, changed representations, targets and noncontent features are checked for every draw. This batch does not repeat maps for the original neural checkpoints.
